\section{Discussion}\label{sec:discussion}

The association between recent unrest and later conflict in the same place is one of
the most robust regularities in event data, and it is often read as escalation. We
have argued that the association has a second component, shared latent episodes that
raise both kinds of events, and that models of conflict alone cannot separate the two
even when they include their own latent space-time field. A joint model in which the
two outcomes share a space-time field can. In simulations it recovers the escalation
coefficient where univariate models report escalation that does not exist, and in
ACLED data for West Africa and the Sahel and for Africa and the Middle East it
attributes between a quarter and two thirds of the univariate association to shared
episodes. Once both shared and unrest-specific episodes are modelled, a short-run
effect of recent unrest on conflict remains, of about two thirds of the univariate
estimate. \todo{Update with the Africa and Middle East refit without Europe (job 29).}

\paragraph{What the shared field represents.} The shared field is a latent component
of \emph{recorded} risk. It combines common shocks to the occurrence of conflict and
unrest, such as a political crisis, with common shocks to their reporting, such as a
period of media attention or a change in the sources that cover a region
\citep{weidmann2016closer, hellmeier2018spotlight, croicu2022reporting}. With a single
source of event data the two cannot be separated, and we do not interpret the field,
or its loading, as a property of occurrence alone. Several independent sources of
event data, or auxiliary measures of reporting intensity, would allow the two to be
distinguished \citep{cook2017twowrongs}. The same applies to the escalation
coefficient: it is a conditional association of recorded events, not a causal
effect.

\paragraph{The scale of escalation.} The joint model separates escalation from shared
episodes only at the resolution of the shared field. Escalation that unfolds over
several months and over hundreds of kilometres is indistinguishable from a shared
episode and is attributed to the field, while episodes that are shorter than a
time step or smaller than the lattice of the field remain in the lag coefficient. Our
estimate is therefore the short-run, local part of escalation. It is compatible with
accounts in which repression of protest leads to armed conflict over longer periods
\citep{ritter2016preventing, ryckman2020turn, chenoweth2017violent, rod2023protest},
and it indicates that four-week unrest histories in the same cell are a weak signal of
such escalation once shared episodes are modelled. It also explains why unrest
histories help forecasting models of conflict \citep{rod2025predicting}: they are
informative about the current episode, whatever its nature.

\paragraph{Prediction.} A more comprehensive model need not predict better. When both
outcomes are observed everywhere and both histories enter as covariates, the shared
field adds little that the univariate models do not already contain, and in our
simulations and data the joint model forecasts neither outcome better. Its predictive
value appears when one outcome is missing where the other is recorded: there the
shared field transfers information about the current episode from the observed
outcome to the unobserved one. Under-reporting of particular types of events is
common in event data \citep{weidmann2016closer, croicu2022reporting}, so this is a
practical use of the model, for example to flag areas where recorded protest
suggests that conflict may be going unrecorded. The gain fades in forecasts, because
the shared field decays at the rate of its autoregression and a shared static
field retained only a small part of it. The joint model also under-predicts the
co-occurrence of the two outcomes in the test period, by about 13--15\% in West
Africa even with a shared static field.

\paragraph{Limitations.} Several limits apply. The plug-in standard errors ignore the
uncertainty in the hyperparameters and, in simulations, understate the variability of
the lag coefficients by up to 20\% when a space-time field is fitted, so intervals
that only just exclude zero should be read with caution. The loading of the shared
field is constant in space, while the relation between unrest and conflict surely
differs between, say, the central Sahel and North Africa. Escalation runs in both
directions in the model, but the simulations covered one direction only, and they
did not reproduce the 13- and 26-week steps or the coarser lattice of the application.
The regions were chosen as boxes and the countries of Europe were excluded from the
larger one; results for other regions may differ. Finally, the analysis uses a
repaired version of an earlier cell-week table, whose construction we describe in the
Supplementary Material.

\paragraph{Extensions.} Natural extensions are loadings that vary in space, through a
second field in the loading or region-specific loadings; non-separable space-time
fields \citep{lindgren2024diffusion}; more outcomes, such as fatalities or different
actors, sharing several latent components; and hybrids with self-exciting models of
event data \citep{jun2024hawkes, browning2026hawkes}, in which the shared field would
act as a common stochastic background rate. The finite-volume construction on the
sphere and the case-control subsampling of non-events make such models feasible on
continental grids with millions of cell-weeks.
